Nguồn: \href{https://vietnamnet.vn/vn/giao-duc/nguoi-thay/gs-phan-dinh-dieu-qua-hoi-uc-cua-con-gai-710146.html}{Bài thơ theo suốt cuộc đời nhà toán học Phan Thị Hà Dương - VietnamNet} và \href{https://vietnamnet.vn/vn/giao-duc/nguoi-thay/gs-phan-dinh-dieu-day-con-710149.html}{Bài học sâu sắc mà GS Phan Đình Diệu dạy con gái - VietnamNet}

\paragraph{Bài thơ theo suốt cuộc đời nhà toán học Phan Thị Hà Dương - VietnamNet} \ \\

"... Nhưng tôi nhớ, chẳng bao giờ tôi có thể quên, bài thơ đầu tiên mà tôi tự cầm sách đọc, đọc và yêu thích, đọc và ghi nhớ, đọc và mang theo suốt cuộc đời. Đó là "Buổi sơ khai"...."

LTS: "Màu sắc của bìa sách, cả việc cuốn sách được bọc và nét chữ của bố như chìm vào bìa mang đến cho tôi một cảm giác trân trọng trang nghiêm và mênh mang" - Ấn tượng sâu đậm nhất về người bố - GS. Phan Đình Diệu - nhà toán học và là người đặt nền móng cho ngành Tin học của Việt Nam, trong hồi ức của người con - cũng là một người làm toán - hoá ra lại không phải về toán mà là... văn. Văn chương, và đặc biệt thơ ca, nơi những cảm xúc trái ngược quyện lại hài hòa, nơi điều gần gụi và cái vô cùng ở sát bên nhau, đã hình thành nên con người và gắn kết sâu sắc thêm tình cha con. 

Báo VietNamNet xin trân trọng giới thiệu hồi ức của nhà toán học Phan Thị Hà Dương về người bố đáng kính của mình.

\paragraph{Bài thơ theo suốt cuộc đời nhà toán học Phan Thị Hà Dương} \ \\

Tôi không còn nhớ bài thơ đầu tiên tôi được học là bài nào, có lẽ đó sẽ là một bài từ thời mẫu giáo, hay xa hơn nữa từ thời nhà trẻ, vì tôi, như tất cả bọn trẻ con thời ấy, đều đi nhà trẻ từ ngày một tuổi.

Tôi không còn nhớ lần đầu tiên tôi biết đọc là khi nào, có thể là trước khi đến trường, vì tôi, như tất cả bọn trẻ con muôn đời, thích tập đọc những bảng chỉ đường, những biển hiệu dọc phố.

Nhưng tôi nhớ, chẳng bao giờ tôi có thể quên, bài thơ đầu tiên mà tôi tự cầm sách đọc, đọc và yêu thích, đọc và ghi nhớ, đọc và mang theo suốt cuộc đời. Đó là "Buổi sơ khai".

Ngày ấy, tôi, con bé con 8 tuổi, mùa hè mặc váy trắng, nằm trên sàn nhà lát những viên gạch đỏ trong phòng ngoài một căn hộ nhỏ trên tầng 4 khu Giảng Võ. Bố tôi cầm một cuốn sách, hình như bìa màu trắng bạc, không có hình vẽ gì, chỉ một dòng chữ nhỏ: “Thơ Tagore”. Bố tôi đưa tôi và chỉ vào một bài thơ: ”Bài này hay lắm, con đọc đi".

Tôi đã đọc bài thơ trong những hình ảnh thiêng liêng về một nhà thơ râu tóc bạc phơ, hiền từ và đẹp như một vị thánh; về sự ngưỡng mộ sùng kính của những người đọc thơ, khi đến thăm ông lúc về đã đi giật lùi. Những câu chuyện bố kể như in vào đầu tôi, con bé con 8 tuổi, về sự linh thiêng và cao quý của Nhà Thơ.

Tôi đã đọc bài thơ, và ngay những dòng đầu đã thích chí cười với “mẹ ôm chặt bé vào lòng và trả lời nửa cười nửa khóc”. Tôi còn nhớ như in cảm giác sung sướng của đứa bé là tôi lúc ấy khi thấy hình ảnh vừa cười vừa khóc, khi cảm thấy mình hiểu được sự mâu thuẫn của hình ảnh vừa mừng vui, vừa ngộ nghĩnh, vừa xúc động.

Tôi còn nhớ tôi đã hỏi mẹ “ban thờ” là gì và khi mẹ trả lời tôi thấy thật hay biết bao khi trong thơ có những từ thật là lạ, và “ban thờ” là một điều gì đó cao xa hơn nhiều những từ mà tôi biết, “ban thờ” chứ không phải “bàn thờ”.

Tôi còn nhớ “và trong cả cuộc đời của mẹ mẹ nữa kia” đã gieo vào tôi ý niệm về một cái gì đó lồng vào nhau, mẹ của mẹ, mẹ của mẹ của mẹ, và xa nữa, xa xa nữa…

Tôi còn nhớ, mỗi lần tôi đọc bài thơ này cho bố mẹ, cho cả nhà nghe, đến đoạn “khi đương thời con gái”, thể nào mẹ tôi cũng đọc cùng tôi, và gương mặt mẹ bừng lên rạng rỡ.

Tôi đã đọc, tôi đã đọc nhiều lắm, tôi không biết vì sao mình lại thích bài thơ này đến vậy. Lúc nào tôi cũng đọc. Họp mặt khu tập thể, tôi đứng lên ghế đọc cho các cô các bác, các bà nghe. Ngày thi văn nghệ theo các khu nhà các phường, các huyện, tôi diện áo trắng juyp xanh xếp nếp đứng trên bục sân khấu đọc cho cả hội trường nghe. Khi ấy, bọn con trai lớp tôi, những cậu chàng lớp 8, khúc khích cười trêu lúc tôi đặt hai tay lên ngực và đọc “mẹ đã siết chặt con trên ngực mẹ”.


Tôi đã đọc những đêm ru con ngủ những ngày đầu làm mẹ. Dẫu rằng bài thơ trúc trắc nhường này nhưng với tôi nó vẫn êm dịu xiết bao. Khi ấy, mẹ tôi, giờ đây đã là mẹ của mẹ, có lẽ lại như ngày xưa nhẩm cùng tôi đoaanj "khi đương thời con gái".

Có thể là ngày bé tôi đã chỉ hiểu lơ mơ. Có thể là với thời gian, tôi đã thấm hiểu niềm hạnh phúc trong câu trả lời nửa cười nửa khóc. Có thể là đến bây giờ, khi tôi đã có hai con, tôi vẫn chưa hiểu hết bài thơ.

Nhưng tôi luôn thầm cám ơn bố mẹ đã cho tôi đọc bài thơ này vào buổi sơ khai của cuộc sống tâm hồn. Để tôi, trong dấu ấn đầu tiên của tâm hồn mình, trong những bước tiếp theo và mãi mãi, luôn mơ màng rằng trong thơ có ẩn chứa những từ ngữ cao xa, những từ ngữ đa nghĩa, rằng trong thơ những cảm xúc trái ngược, những cảm xúc mâu thuẫn bỗng quyện lại hài hòa, rằng trong thơ điều gần gụi và cái vô cùng sát bên nhau đến thế. Và nếu như nhà thơ viết một bài thơ không chỉ bằng một phút giây tỏa sáng mà bằng một phút giây tỏa sáng cộng với cả cuộc đời mình; thì người đọc thơ đọc một bài thơ không chỉ bằng một đêm xuân khi hình như mưa lất phất bay mà bằng một đêm xuân mưa lất phất bay cộng với cả cuộc đời mình.

Buổi sơ khai

Bé hỏi mẹ:

“Mẹ ơi, con từ đâu đến vậy

Mẹ đã nhặt được con ở tận nơi nào ?”

Mẹ ôm chặt bé vào lòng, và trả lời

nửa cười nửa khóc:

“Con ơi con, con đã được giấu kín trong lòng mẹ

như chính những thèm khát ước mơ của nó.

Con ở trong con búp bê của những món đồ chơi tuổi nhỏ của mẹ.

Và mỗi buổi sáng khi mẹ lấy đất sét nặn ra

hình ảnh Chúa Đời của mẹ

thì mẹ đã nặn đi nặn lại con rồi

Con ở trên ban thờ nơi thờ vị thổ thần

Và khi thờ thần đó, đồng thời mẹ cũng thờ con

Con đã sống trong tất cả mọi niềm hy vọng, thương yêu trong đời mẹ

và trong cuộc đời của mẹ mẹ nữa kia

Con đã được nuôi dưỡng đời này qua đời khác

trong lòng của vị thần linh bất tử đã ngự trị ở nhà ta.

Khi đương thời con gái, trái tim mẹ nở xòe như một đóa hoa.

Con đã lượn quanh nó như một mùi hương phảng phất

Vẻ tươi mát nhẹ nhàng của con

Nở trên chân tay non trẻ của mẹ

như một ánh hồng

trên trời cao

trước buổi bình minh

Con là đứa con cưng của Thượng đế,

là anh em sinh đôi với ánh bình minh

Con đã theo dòng nước trôi xuống cuộc đời trần tục này

và cuối cùng con đã được đặt vào trong lòng mẹ

Khi mẹ ngây nhìn khuôn mặt của con

mẹ như bị ngập trong bao điều bí ẩn,

Và con, vốn là của chung của tất cả mọi người

đã trở thành của riêng của mẹ

Sợ mất con đi, mẹ đã siết chặt con trên ngực mẹ

Không hiểu sự diệu kỳ nào

Đã chiếm lĩnh cái kho vàng trên cõi thế

và đặt con vào đôi tay mảnh khảnh của mẹ đây ?”

(Thơ Tagore, Đào Xuân Quý dịch) 

\paragraph{Bài học sâu sắc mà GS Phan Đình Diệu dạy con gái - VietnamNet}

Trong ký ức của nhà toán học Phan Thị Hà Dương, chỉ có một lần duy nhất GS Phan Đình Diệu - can thiệp vào việc học văn của con gái: “Bố muốn con hiểu rằng con chỉ nên viết ra những gì mà con thực sự nghĩ là đúng”.

Thời phổ thông, có nhiều lúc tôi ngán mớ đời khi có đứa bạn hay thậm chí thầy giáo hỏi có phải cái bài toán về nhà tôi nhờ bố nên mới giải được, trong khi thực ra toàn bộ đời đi học của tôi chỉ có hai bài nhờ bố giải. Ngay cả khi tôi hỏi bố nên học toán hay tin học thì bố cũng bảo: “Bố tôn trọng mọi quyết định của con”.

Vậy mà trong thâm tâm tôi luôn nghĩ bố là người thầy dạy mình học văn dù rằng chỉ có một lần duy nhất bố can thiệp vào việc học văn của tôi.

Lần ấy tôi bị 4 điểm bài văn kiểm tra chất lượng đầu năm lớp 6. Trước đấy, mặc dù cứ thích đọc thơ và suốt ngày ôm truyện nhưng đi học thì tôi chỉ khoái môn toán nên cũng chẳng ghét bỏ hay yêu quý gì môn văn, có khi tôi được 9 điểm cao nhất hồi lớp 3 khi vô cùng xúc động làm bài văn về anh Lê Văn Tám, cũng có khi tôi làm lục bát ngang phè vào hồi lớp 5. Nói chung không có vấn đề gì nghiêm trọng.

Nhưng 4 điểm, lần đầu tiên dưới trung bình, là chuyện tày đình, và vì thế lần đầu tiên bố đọc bài văn của tôi. Tôi vẫn nhớ có một câu đề bài là “Hãy viết một câu văn hay, sử dụng biện pháp nhân cách hóa” và tôi đã chép nguyên xi một trong các câu mẫu của cô “những dãy nhà tường trắng mái đỏ của trường em trông như những chú lùn đội mũ đỏ”. 

Bố bảo tôi: “Bố không thấy câu văn này hay ở chỗ nào”. Tôi lý sự: “Nhưng chính cô cho bọn con học câu mẫu, bây giờ con chép lại thì cô lại trừ điểm”. Bố nhìn tôi rất nghiêm trang và hỏi: “Thế con có thấy những dãy nhà ở trường giống các chú lùn không?”. Tôi lí nhí: “Nhưng mà cô ...”. "Không, bố hỏi con cơ, con có thấy thế không?”. Tôi được dịp bùng ra: “Không ạ, con chẳng thấy giống gì cả, nhà thì dài dằng dặc, như cái hộp, chẳng biết sao cô lại cho như thế, chỉ tại cái nhân cách hóa thôi ạ!".

Lúc ấy bố nhìn tôi rất thẳng và bảo tôi: “Bố muốn con hiểu rằng con chỉ nên viết ra những gì mà con thực sự nghĩ là đúng.”

Còn tôi làm nũng bố: “Con chẳng biết thế nào là một câu văn hay cả”.  

Bố nói rằng sao bố thấy con suốt ngày đọc truyện mà bây giờ đến một câu văn hay cũng không biết, con đang đọc truyện gì vậy. Thật là xui xẻo cho tôi, cái thời đấy thì vớ được gì đọc nấy chứ có nhiều lựa chọn đâu (hồi ấy phải có người quen mới mua được truyện từ NXB mà). Nếu mà cách đấy mươi ngày, tôi đang đọc “Tôm Giôn- đứa trẻ vô thừa nhận” tập 1 (tập 2 và 3 phải cả năm sau mới được đọc) thì đã chẳng vấn đề gì, đằng này lúc ấy tôi lại đang một gối hai quyển dày cộm “Ghenny Ghéchac” và “Hoa hậu xứ Mường”. Bố không bằng lòng một chút nào, bố bảo rằng xưa nay bố để tôi tự do đọc sách nhưng vì bây giờ tôi không biết viết một câu văn nữa nên bố mới phải quan tâm đến việc này, và rằng hai cuốn này không phù hợp với tôi. Dù tôi đã năn nỉ bố là đang đúng đoạn hồi hộp và truyện sắp phải trả rồi nhưng bố kiên quyết không.

Rồi bố bảo: “Nếu con muốn biết thế nào là một câu văn hay thì con có thể đọc cuốn này”, và bố rút trên giá sách xuống một cuốn sách. Đến bây giờ tôi vẫn nhớ cảm giác khi cầm cuốn sách đó. Một cuốn sách khá mỏng, bìa được bọc bằng giấy nâu, kiểu giấy xi măng ngày xưa, ở gáy sách và ở bìa sách, phía chếch lên bên phải là những nét chữ in của bố, nét chữ mực màu đen: “GIÓ ĐẦU MÙA”. Màu sắc của bìa sách, cả việc cuốn sách được bọc và nét chữ của bố như chìm vào bìa mang đến cho tôi môt cảm giác trân trọng trang nghiêm và mênh mang.

Tôi vẫn nhớ hai truyện đầu tiên là “Nhà mẹ Lê” và “Hai đứa trẻ”, chỉ có điều tôi không nhớ là truyện nào trước truyện nào sau (có lần tôi tranh luận với một bạn, tôi nói “Nhà mẹ Lê mình nhớ rõ ràng” còn nó bảo “Hai đứa trẻ không thể sai được”, cuối cùng cùng với thời gian thì ý kiến của nó hòa vào với tôi đến nỗi bây giờ tôi mới chẳng còn nhớ gì thế này).

Sau này, theo trí nhớ của tôi, tôi đã luôn tự viết các bài văn mà chẳng bao giờ chép lại từ đâu cả, và có vẻ là tôi còn rất tự tin nữa cơ.

Sau Thạch Lam, tôi quá háo hức, và đã quét sạch cả một ngăn giá sách của bố, đầu tiên là mấy tập tuyển tập Nam Cao, tôi thấy buồn quá và cũng nhiều chuyện không hiểu được, rồi thích nhất là hai (hay ba) tập Nguyễn Công Hoan, văn ông rất sinh động, nhiều đối thoại; rồi Thế Lữ. Thế Lữ cả thơ và văn chỉ một quyển, rất dày. Phần đầu là thơ, mở đầu bằng “Nhớ rừng”, phần sau là văn. Tôi kể cho bố nghe tôi khoái chí thế nào khi đọc “Những nét chữ” và “Vàng và máu”, bố bảo hồi sinh viên bố cũng rất thích đọc truyện trinh thám của Thế Lữ. Bố mẹ và cả nhà rất thích nghe tôi đọc “Nhớ rừng”, tôi bé con mà rất thích đọc hùng tráng, còn bố thích đọc “Tiếng trúc tuyệt vời”, tôi vẫn nhớ những lúc bố đọc:

“Cô em đứng bên hồ

nghiêng tựa mình cây, dáng thẩn thơ

Chừng cô tưởng đến ngày vui sẽ mất,

mà sắc đẹp rỡ ràng rồi sẽ tắt

Cho nên khi cô nghe tiếng trúc tuyệt vời

Thổn thức với lòng cô thổn thức

Man mác với lòng cô man mác

Cô để tâm hồn tê tái bâng khuâng ..."

Và bố hỏi tôi có biết vì sao lại gọi là “cô em” mà không phải là “cô” hay “em” không, nhưng hình như bố không trả lời. Sau này, khi đọc Thi nhân Việt Nam tôi có thấy Hoài Thanh bình chữ này.

Ngày ấy, đi đâu bố cũng xách tôi theo, tôi đi gặp những ông Moi Sép, cô Linđa, tôi lên Đồi Thông, tôi đến Nghĩa Đô.

Hồi tôi mới ở Pháp về Viện Toán làm việc, có lần tôi đang đứng ở bảng thông báo đọc linh tinh thì bỗng nghe: “Em có phải là Hà Dương không, sao em lại ở đây, trông em chẳng khác cái hồi 9 tuổi hay lên Viện đọc thơ gì cả”. Tôi buồn cười quá, vì tôi nghĩ là mình rất khác, thậm chí khác với vài tháng trước đó. Và mặc dù tất nhiên là tôi chẳng nhớ ra anh ấy là ai nhưng tôi vẫn được tặng một bó hoa bất tử đang bày trong phòng làm việc của anh ấy. Bó hoa ấy tôi vẫn để trong phòng làm việc của mình (hơi bụi một tí) như kỷ niệm về một thời 9 tuổi hay lẽo đẽo theo chân bố mẹ.

Một chấm sao không ngủ cuối thiên hà

Ngày trước, những năm 80, những năm cuối cấp một và cấp hai của tôi, nhà tôi ở khu Đồng Xa, bố tôi có rất nhiều bạn bè đến chơi. Chủ nhật nào phòng khách nhà tôi cũng có bạn của bố tôi, những bạn học cũ, bạn toán, bạn văn; những người bạn mới, có những người vì một bài viết của bố tôi mà đã đến rất nhiều chủ nhật và đã trở nên thân thiết.

Tôi vẫn nhớ chiều chủ nhật ấy, trong phòng khách nhà tôi có bố mẹ tôi, có cậu Cương (PGS. Văn Như Cương), có bác Đoàn Quỳnh, bác Hoàng Xuân Sính, có lẽ có cả bác Hà Văn Tấn nữa, và các bác khác. Bố tôi nói: "Mình có tập thơ này hay lắm, của Việt Phương”. Và bố tôi giới thiệu với mọi người một tập bản thảo chép tay của bác Việt Phương, nét bút mực trên nền giấy hơi ngà sẫm màu, các chữ đầu dòng đều không viết hoa, và tên bài thơ nào cũng chỉ là một chữ. Tập thơ ấy không phổ biến và bác đã cho bố mượn.

Bố tôi đọc cho các bạn mình nghe một số bài. Tôi vẫn nhớ không khí của buổi chiều ấy, niềm hứng khởi và sự tâm đắc của bố và các bạn. Theo trí nhớ của tôi thì nhiều bài thơ mang tính trí tuệ và mọi người đã thán phục vì những tứ thơ độc đáo và sâu cay, có những tứ thơ làm mọi người bật lên như một sự khám phá. Nhưng bài thơ mà tôi thích nhất là một bài thơ tình cảm, với một cái tên thật lạ: màu. Tôi vẫn nhớ giọng đọc thơ của bố khi đó, tách khỏi giọng đọc có phần nhấn nhá, đôi khi hơi hài hước và có khi nhấn giọng lúc trước, bài thơ này bố tôi đọc rất tình cảm.

Đến bây giờ tôi vẫn như đang nghe thấy giọng đọc trầm ấm, trìu mến và tình cảm của bố tôi.

em cứ là những tinh mơ tê tái rét

phanh cổ áo ra cho gió siết vào da

em cứ là cơn giông đầu mùa

đi đầu trần đón dòng mưa xối xả

em cứ là cái khoảng cách chập chờn sương phủ

suốt một đời anh vất vả vượt qua

em cứ là giữa mịt mùng vô định

một chấm sao không ngủ cuối thiên hà  

Những tối sau đó, bố tôi còn đọc cho mấy mẹ con nghe, có những bài bố cho tôi đọc nữa. Chỉ hơn một tuần thôi, rồi bố tôi đã trả lại tập bản thảo cho bác Việt Phương. Nhưng bài thơ đã in vào trong trí não tôi.

Sau này, sau những năm 1990, khi nhà tôi đã chuyển, khi những cuộc cách mạng đã nở bừng trên thế giới, khi những biến cố lớn đã đến với biết bao người bạn thân thiết của bố tôi, người ta đã nhắc nhiều hơn đến những bài thơ của bác. Và mãi về sau, khi tập thơ “Cửa đã mở" của bác Việt Phương được in, tôi đã tìm mua, mong nhìn lại những bài thơ hồi bé tôi đã được nghe bố đọc. Tôi tìm thấy những bài như bài “Thịt":

Chị mười ba ý tứ nết na

Cuối bữa cơm gắp rụt rè một miếng

(Ngày trước, khi bố tôi đọc bài này, mẹ tôi hay bảo giống chị tôi, lúc nào cũng nhường nhịn cả nhà).

Có nhiều bài nữa, nhưng tôi không tìm thấy bài thơ trong tâm trí tôi. Và vì thế, có những khi tôi cứ thưởng cho mình cái ý nghĩ rằng chỉ có bố tôi và tôi nhớ bài thơ ấy thôi.

Đến bây giờ tôi vẫn chưa hiểu rõ hết nghĩa của tên bài thơ, nhưng rất nhiều những sáng tinh mơ khi tiết thu đã hết và những làn gió sớm mùa đông siết buốt da, tôi lại phóng trên đường phố Hà Nội, lại bỏ khăn quàng cổ để cảm nhận câu thơ.

Đến bây giờ, khi viết những dòng này, tôi chợt nhận ra vì sao việc đọc một bài thơ đối với tôi có ý nghĩa thiêng liêng đến thế. Tôi đã chịu ảnh hưởng của bố, đã luôn nâng niu từng bài thơ, nâng niu từng giây phút mình đọc thơ. Tôi đã luôn chịu ảnh hưởng của bố, từ ngày bé thơ cho đến sau này, và mãi mãi.

"Và nếu như nhà thơ viết một bài thơ không chỉ bằng một phút giây tỏa sáng mà bằng một phút giây tỏa sáng cộng với cả cuộc đời mình; thì người đọc thơ đọc một bài thơ không chỉ bằng một đêm xuân khi hình như mưa lất phất bay mà bằng một đêm xuân mưa lất phất bay cộng với cả cuộc đời mình...".